%%%%%%%%%%%%%%%%
%%% num 14 %%%
%%%%%%%%%%%%%%%%

\Chapter{14}

\Verse{1}
{
	\hRJE{וַתִּשָּׂא כּל הָעֵדָה וַיִּתְּנוּ אֶת קוֹלָם וַיִּבְכּוּ הָעָם בַּלַּיְלָה הַהוּא׃}
}
{
	\eRJE{
		And all the congregation raised and let out their
		voices! And the people wept that night.
	}
}
{
}

\Verse{2}
{
	\hRJE{וַיִּלֹּנוּ עַל מֹשֶׁה וְעַל אַהֲרֹן כֹּל בְּנֵי יִשְׂרָאֵל וַיֹּאמְרוּ אֲלֵהֶם כּל הָעֵדָה לוּ מַתְנוּ בְּאֶרֶץ מִצְרַיִם אוֹ בַּמִּדְבָּר הַזֶּה לוּ מָתְנוּ׃}
}
{
	\eRJE{
		And all the children of Israel complained at Moses and at
		Aaron, and all the congregation said to them, “If only we
		had died in the land of Egypt! Or in this wilderness, if
		only we had died!
	}
}
{
}

\Verse{3}
{
	\hRJE{וְלָמָה יְהֹוָה מֵבִיא אֹתָנוּ אֶל הָאָרֶץ הַזֹּאת לִנְפֹּל בַּחֶרֶב נָשֵׁינוּ וְטַפֵּנוּ יִהְיוּ לָבַז הֲלוֹא טוֹב לָנוּ שׁוּב מִצְרָיְמָה׃}
}
{
	\eRJE{
		And why is YHWH bringing us to this land to fall by the
		sword? Our wives and our infants will become a spoil!
		Isn’t it better for us to go back to Egypt?”
	}
}
{
}

\Verse{4}
{
	\hRJE{וַיֹּאמְרוּ אִישׁ אֶל אָחִיו נִתְּנָה רֹאשׁ וְנָשׁוּבָה מִצְרָיְמָה׃}
}
{
	\eRJE{
		And they said, each man to his brother, “Let’s appoint a
		chief and go back to Egypt.”
	}
}
{
}

\Verse{5}
{
	\hRJE{וַיִּפֹּל מֹשֶׁה וְאַהֲרֹן עַל פְּנֵיהֶם לִפְנֵי כּל קְהַל עֲדַת בְּנֵי יִשְׂרָאֵל׃}
}
{
	\eRJE{
		And Moses and Aaron fell on their faces in front of all
		the community of the congregation of the children of
		Israel.
	}
}
{
}

\Verse{6}
{
	\hRJE{וִיהוֹשֻׁעַ בִּן נוּן וְכָלֵב בֶּן יְפֻנֶּה מִן הַתָּרִים אֶת הָאָרֶץ קָרְעוּ בִּגְדֵיהֶם׃}
}
{
	\eRJE{
		And Joshua son of Nun and Caleb son of Jephunneh, from
		those who had scouted the land, had torn their clothes.
	}
}
{
}

\Verse{7}
{
	\hRJE{וַיֹּאמְרוּ אֶל כּל עֲדַת בְּנֵי יִשְׂרָאֵל לֵאמֹר הָאָרֶץ אֲשֶׁר עָבַרְנוּ בָהּ לָתוּר אֹתָהּ טוֹבָה הָאָרֶץ מְאֹד מְאֹד׃}
}
{
	\eRJE{
		And they said to all the congregation of the children of
		Israel, saying, “The land through which we passed to scout
		it: the land is very, very good!
	}
}
{
}

\Verse{8}
{
	\hRJE{אִם חָפֵץ בָּנוּ יְהֹוָה וְהֵבִיא אֹתָנוּ אֶל הָאָרֶץ הַזֹּאת וּנְתָנָהּ לָנוּ אֶרֶץ אֲשֶׁר הִוא זָבַת חָלָב וּדְבָשׁ׃}
}
{
	\eRJE{
		If YHWH desires us, then He’ll bring us to this land and
		give it to us, a land that flows with milk and honey!
	}
}
{
}

\Verse{9}
{
	\hRJE{אַךְ בַּיהֹוָה אַל תִּמְרֹדוּ וְאַתֶּם אַל תִּירְאוּ אֶת עַם הָאָרֶץ כִּי לַחְמֵנוּ הֵם סָר צִלָּם מֵעֲלֵיהֶם וַיהֹוָה אִתָּנוּ אַל תִּירָאֻם׃}
}
{
	\eRJE{
		Just don’t revolt against YHWH. And you, don’t fear the
		people of the land, because they’re our bread! Their
		protection has turned from them, while YHWH is with us.
		Don’t fear them.”
	}
}
{
}

\Verse{10}
{
	\hRJE{וַיֹּאמְרוּ כּל הָעֵדָה לִרְגּוֹם אֹתָם בָּאֲבָנִים וּכְבוֹד יְהֹוָה נִרְאָה בְּאֹהֶל מוֹעֵד אֶל כּל בְּנֵי יִשְׂרָאֵל׃}
}
{
	\eRJE{
		And all the congregation said to batter them with stones.
		And YHWH’s glory appeared at the Tent of Meeting to all
		the children of Israel!
	}
}
{
}

\Verse{11}
{
	\hJ{וַיֹּאמֶר יְהֹוָה אֶל מֹשֶׁה עַד אָנָה יְנַאֲצֻנִי הָעָם הַזֶּה וְעַד אָנָה לֹא יַאֲמִינוּ בִי בְּכֹל הָאֹתוֹת אֲשֶׁר עָשִׂיתִי בְּקִרְבּוֹ׃}
}
{
	\eJ{
		And YHWH said to Moses, “How long will this people reject
		me, and how long will they not trust in me, with all the
		signs that I’ve done among them?
	}
}
{
%	\footnote{
%		how long will they not trust in me, with all the signs that I’ve done
%		among them? Why in fact do the people not trust? Why are they still
%		afraid and doubtful? They have seen their God handle Egypt, the major
%		power of the world. Why would they doubt that He can handle the
%		Canaanites? They have seen Him manipulate the whole of nature with the
%		ten plagues and the splitting of the sea. They have seen water come from
%		a mountain crag. They have heard God’s voice from the sky at Sinai. They
%		are eating miraculous manna every day. They have just miraculously been
%		provided with quail to eat. Even as they complain in this moment, there
%		is a miraculous column of cloud and fire right in front of them. Why are
%		they afraid? Recall that the reason they are directed to the Red Sea
%		when they leave Egypt, rather than through Philistine territory, is that
%		YHWH is concerned that they will be afraid of facing war with the
%		Philistines. As I commented there (Exod 13:18), the Red Sea required no
%		decisions or action on their part, but facing the Philistines would have
%		required them to take up arms. Likewise, in every case so far, from the
%		plagues to the quail, God and Moses have acted for them, and they have
%		had to do very little. Even in the battle against the Amalekites, it was
%		the Amalekites who attacked them, and then Joshua chose a select group
%		to fight on their behalf (Exod 17:8–13). But now, for the first time,
%		the people perceive that they will have to act. The scouts speak in
%		terms of the people of Israel themselves having to fight: “We won’t be
%		able to go up against the people, because they’re stronger than we are”
%		(13:31). The scouts do not even mention God. What was true when they
%		left Egypt is still true: They are still slaves in their hearts. They
%		are not yet capable of independence, responsibility, or action. And so
%		God’s decision is precisely to leave the entry into the land for their
%		children. See the comment on Num 14:29.
%	}
}

\Verse{12}
{
	\hJ{אַכֶּנּוּ בַדֶּבֶר וְאוֹרִשֶׁנּוּ וְאֶעֱשֶׂה אֹתְךָ לְגוֹי גָּדוֹל וְעָצוּם מִמֶּנּוּ׃}
}
{
	\eJ{
		I’ll strike them with an epidemic and dispossess them, and
		I’ll make you into a bigger and more powerful nation than
		they are.”
	}
}
{
%	\footnote{
%		I’ll make you. The Septuagint and Samaritan texts have “I’ll make you
%		and your father’s house.” The extra words appear to have been lost by
%		haplography. (That is, in the Hebrew, the word for “make you” and the
%		phrase for “your father’s house” both end in the same letter, kaf, so a
%		scribe’s eye skipped from one kap to the next and left out the words in
%		between.) The addition of the reference to Moses’ father’s house may
%		suggest that Aaron’s family is to be spared as well, or it may mean more
%		broadly that the Levites as a whole are to survive. Alternatively, it
%		may still mean that only Moses is to survive to start over a new people,
%		but it is still described as the survival of his father’s house because
%		God’s relationship with Moses has focused on that link from the
%		beginning—God’s self-introduction to Moses at the bush is: “I am your
%		father’s God” (Exod 3:6).
%	}
}

\Verse{13}
{
	\hJ{וַיֹּאמֶר מֹשֶׁה אֶל יְהֹוָה וְשָׁמְעוּ מִצְרַיִם כִּי הֶעֱלִיתָ בְכֹחֲךָ אֶת הָעָם הַזֶּה מִקִּרְבּוֹ׃}
}
{
	\eJ{
		And Moses said to YHWH, “And Egypt will hear it, for you
		brought this people up from among them with your power,
	}
}
{
}

\Verse{14}
{
	\hJ{וְאָמְרוּ אֶל יוֹשֵׁב הָאָרֶץ הַזֹּאת שָׁמְעוּ כִּי אַתָּה יְהֹוָה בְּקֶרֶב הָעָם הַזֶּה אֲשֶׁר עַיִן בְּעַיִן נִרְאָה אַתָּה יְהֹוָה וַעֲנָנְךָ עֹמֵד עֲלֵהֶם וּבְעַמֻּד עָנָן אַתָּה הֹלֵךְ לִפְנֵיהֶם יוֹמָם וּבְעַמּוּד אֵשׁ לָיְלָה׃}
}
{
	\eJ{
		and they’ll say it to those who live in this land. They’ve
		heard that you, YHWH, are among this people; that you,
		YHWH, have appeared eye-to-eye; and your cloud stands over
		them; and you go in front of them in a column of cloud by
		day and in a column of fire by night.
	}
}
{
}

\Verse{15}
{
	\hJ{וְהֵמַתָּה אֶת הָעָם הַזֶּה כְּאִישׁ אֶחָד וְאָמְרוּ הַגּוֹיִם אֲשֶׁר שָׁמְעוּ אֶת שִׁמְעֲךָ לֵאמֹר׃}
}
{
	\eJ{
		And if you kill this people as one man, then the nations
		that have heard about you will say it, saying,
	}
}
{
}

\Verse{16}
{
	\hJ{מִבִּלְתִּי יְכֹלֶת יְהֹוָה לְהָבִיא אֶת הָעָם הַזֶּה אֶל הָאָרֶץ אֲשֶׁר נִשְׁבַּע לָהֶם וַיִּשְׁחָטֵם בַּמִּדְבָּר׃}
}
{
	\eJ{
		‘Because YHWH wasn’t able to bring this people to the land
		that He swore to them, He slaughtered them in the
		wilderness.’
	}
}
{
}

\Verse{17}
{
	\hJ{וְעַתָּה יִגְדַּל נָא כֹּחַ אֲדֹנָי כַּאֲשֶׁר דִּבַּרְתָּ לֵאמֹר׃}
}
{
	\eJ{
		And now, let my Lord’s power be big, as you spoke, saying,
	}
}
{
%	\footnote{
%		let my Lord’s power be big. Moses suggests that it is more powerful to
%		be merciful than to punish. In some circumstances, compassion is
%		weakness. In others, it requires enormous strength.
%	}
}

\Verse{18}
{
	\hJ{יְהֹוָה אֶרֶךְ אַפַּיִם וְרַב חֶסֶד נֹשֵׂא עָוֺן וָפָשַׁע וְנַקֵּה לֹא יְנַקֶּה פֹּקֵד עֲוֺן אָבוֹת עַל בָּנִים עַל שִׁלֵּשִׁים וְעַל רִבֵּעִים׃}
}
{
	\eJ{
		‘YHWH is slow to anger and abounding in kindness, bearing
		crime and offense; though not making one innocent:
		reckoning fathers’ crime on children, on third generations
		and on fourth generations.’
	}
}
{
%	\footnote{
%		slow to anger … Moses quotes God’s own words from the moment of the
%		great revelation, when Moses saw God at Sinai (Exod 34:6–7). So he is
%		not only appealing to God’s mercy; he is citing God’s own description of
%		Himself as if to say that God must be true to Himself. Footnote 14:18.
%		on children, on third generations and on fourth generations. The passage
%		that Moses is quoting says: “on children and on children’s children, on
%		third generations and on fourth generations” (Exod 34:7). It looks like
%		“and on children’s children” has been lost, presumably by another
%		haplography: a scribe’s eye skipped from one occurrence of the word
%		“children” to the next. It is the easiest kind of scribal error to make,
%		and in fact this same omission occurs in the Ten Commandments (Exod 20:5
%		and Deut 5:9). As a less banal and more homiletical explanation of this
%		case, we might suggest that the text pictures Moses as shortening the
%		long formula so as to keep focused on the larger point. Or perhaps, by
%		leaving it out, he subtly presses God to make the correction and thus to
%		be reminded of the extent of His compassion according to His own words.
%		Deliberately misquoting someone may actually be an effective tool of
%		persuasion.
%	}
}

\Verse{19}
{
	\hJ{סְלַח נָא לַעֲוֺן הָעָם הַזֶּה כְּגֹדֶל חַסְדֶּךָ וְכַאֲשֶׁר נָשָׂאתָה לָעָם הַזֶּה מִמִּצְרַיִם וְעַד הֵנָּה׃}
}
{
	\eJ{
		Forgive this people’s crime in proportion to the magnitude
		of your kindness and as you’ve borne this people from
		Egypt to here.”
	}
}
{
%	\footnote{
%		Forgive this people’s crime. Moses’ four-fold defense of the people is
%		brilliantly merged: What will the Egyptians say! What will the nations
%		say! You swore the land to them (namely, in the covenant)! You are
%		merciful.
%	}
}

\Verse{20}
{
	\hJ{וַיֹּאמֶר יְהֹוָה סָלַחְתִּי כִּדְבָרֶךָ׃}
}
{
	\eJ{
		And YHWH said, “I’ve forgiven according to your word.
	}
}
{
%	\footnote{
%		I’ve forgiven according to your word. Abraham questioned God regarding
%		Sodom and Gomorrah, but his conversation did not change anything. Ten
%		virtuous persons were not found, and the cities were destroyed. But
%		Moses is pictured as actually making a difference in a divine decision
%		and in the fate of the people. This is another case in which there is
%		growth in the human stance relative to God.
%	}
}

\Verse{21}
{
	\hJ{וְאוּלָם חַי אָנִי וְיִמָּלֵא כְבוֹד יְהֹוָה אֶת כּל הָאָרֶץ׃}
}
{
	\eJ{
		But indeed, as I live, and as YHWH’s glory has filled all
		the earth,
	}
}
{
}

\Verse{22}
{
	\hJ{כִּי כל הָאֲנָשִׁים הָרֹאִים אֶת כְּבֹדִי וְאֶת אֹתֹתַי אֲשֶׁר עָשִׂיתִי בְמִצְרַיִם וּבַמִּדְבָּר וַיְנַסּוּ אֹתִי זֶה עֶשֶׂר פְּעָמִים וְלֹא שָׁמְעוּ בְּקוֹלִי׃}
}
{
	\eJ{
		I swear that all of these people, who have seen my glory
		and my signs that I did in Egypt and in the wilderness and
		who have tested me ten times now and haven’t listened to
		my voice,
	}
}
{
}

\Verse{23}
{
	\hJ{אִם יִרְאוּ אֶת הָאָרֶץ אֲשֶׁר נִשְׁבַּעְתִּי לַאֲבֹתָם וְכל מְנַאֲצַי לֹא יִרְאוּהָ׃}
}
{
	\eJ{
		won’t see the land that I swore to their fathers, and all
		those who rejected me won’t see it.
	}
}
{
}

\Verse{24}
{
	\hJ{וְעַבְדִּי כָלֵב עֵקֶב הָיְתָה רוּחַ אַחֶרֶת עִמּוֹ וַיְמַלֵּא אַחֲרָי וַהֲבִיאֹתִיו אֶל הָאָרֶץ אֲשֶׁר בָּא שָׁמָּה וְזַרְעוֹ יוֹרִשֶׁנָּה׃}
}
{
	\eJ{
		And my servant Caleb, because a different spirit was with
		him, and he went after me completely, I’ll bring him to
		the land where he went, and his seed will possess it.
	}
}
{
}

\Verse{25}
{
	\hJ{וְהָעֲמָלֵקִי וְהַכְּנַעֲנִי יוֹשֵׁב בָּעֵמֶק מָחָר פְּנוּ וּסְעוּ לָכֶם הַמִּדְבָּר דֶּרֶךְ יַם סוּף׃}
}
{
	\eJ{
		And the Amalekite and the Canaanite live in the valley.
		Turn and travel tomorrow to the wilderness by the way of
		the Red Sea.”
	}
}
{
}

\Verse{26}
{
	\hP{וַיְדַבֵּר יְהֹוָה אֶל מֹשֶׁה וְאֶל אַהֲרֹן לֵאמֹר׃}
}
{
	\eP{And YHWH spoke to Moses and to Aaron, saying,}
}
{
}

\Verse{27}
{
	\hP{עַד מָתַי לָעֵדָה הָרָעָה הַזֹּאת אֲשֶׁר הֵמָּה מַלִּינִים עָלָי אֶת תְּלֻנּוֹת בְּנֵי יִשְׂרָאֵל אֲשֶׁר הֵמָּה מַלִּינִים עָלַי שָׁמָעְתִּי׃}
}
{
	\eP{
		“How much further for this bad congregation, that they’re
		complaining against me? I’ve heard the complaints of the
		children of Israel that they’re making against me.
	}
}
{
}

\Verse{28}
{
	\hP{אֱמֹר אֲלֵהֶם חַי אָנִי נְאֻם יְהֹוָה אִם לֹא כַּאֲשֶׁר דִּבַּרְתֶּם בְּאזְנָי כֵּן אֶעֱשֶׂה לָכֶם׃}
}
{
	\eP{
		Say to them: As I live—word of YHWH—what you have spoken
		in my ears, that is what I’ll do to you!
	}
}
{
%	\footnote{
%		what you have spoken in my ears, that is what I’ll do to you! Another
%		case of ironic divine punishment to suit the crime: They could have had
%		the land and lived, but they were afraid and said that they would die
%		and their children would be made captives. So now they will in fact die,
%		but their children will live and one day be free in the land. This is
%		the most dramatic demonstration yet of the people’s persistent slave
%		mentality. God says that He has forgiven them, and they are not stricken
%		with an epidemic as God had originally said; but they are also not ready
%		for the life of free people in a land of their own. And so this is left
%		for their children, who grow up not knowing the experience of having
%		been slaves. An entire generation’s experience is not easily reversed.
%		In our own age we see the impact on the lives of survivors of the
%		holocaust. Some become more religious, some reject religion. Some cannot
%		speak of it, and some speak of it constantly. Some return to the sites
%		where they suffered, and some will never go back. It affects their
%		medical decisions, their ways of raising their children, their trust of
%		other human beings, their work, their thoughts, their opinions, and
%		their memories. Their children carry tremendous loads of guilt and
%		obligation, but they have a chance of knowing a kind of freedom that
%		their parents cannot know.
%	}
}

\Verse{29}
{
	\hP{בַּמִּדְבָּר הַזֶּה יִפְּלוּ פִגְרֵיכֶם וְכל פְּקֻדֵיכֶם לְכל מִסְפַּרְכֶם מִבֶּן עֶשְׂרִים שָׁנָה וָמָעְלָה אֲשֶׁר הֲלִינֹתֶם עָלָי׃}
}
{
	\eP{
		In this wilderness your carcasses will fall; and all of
		you who were counted, for all your number, from twenty
		years old and up, who complained against me,
	}
}
{
}

\Verse{30}
{
	\hP{אִם אַתֶּם תָּבֹאוּ אֶל הָאָרֶץ אֲשֶׁר נָשָׂאתִי אֶת יָדִי לְשַׁכֵּן אֶתְכֶם בָּהּ כִּי אִם כָּלֵב בֶּן יְפֻנֶּה וִיהוֹשֻׁעַ בִּן נוּן׃}
}
{
	\eP{
		I swear that you won’t come to the land that I raised my
		hand to have you reside there—except Caleb son of
		Jephunneh and Joshua son of Nun.
	}
}
{
}

\Verse{31}
{
	\hP{וְטַפְּכֶם אֲשֶׁר אֲמַרְתֶּם לָבַז יִהְיֶה וְהֵבֵיאתִי אֹתָם וְיָדְעוּ אֶת הָאָרֶץ אֲשֶׁר מְאַסְתֶּם בָּהּ׃}
}
{
	\eP{
		And your infants, whom you said would become a spoil: I’ll
		bring them, and they will know the land that you rejected!
	}
}
{
}

\Verse{32}
{
	\hP{וּפִגְרֵיכֶם אַתֶּם יִפְּלוּ בַּמִּדְבָּר הַזֶּה׃}
}
{
	\eP{
		And you: your carcasses will fall in this wilderness.
	}
}
{
}

\Verse{33}
{
	\hP{וּבְנֵיכֶם יִהְיוּ רֹעִים בַּמִּדְבָּר אַרְבָּעִים שָׁנָה וְנָשְׂאוּ אֶת זְנוּתֵיכֶם עַד תֹּם פִּגְרֵיכֶם בַּמִּדְבָּר׃}
}
{
	\eP{
		And your children will be roving in the wilderness forty
		years, and they’ll bear your whoring until the end of your
		carcasses in the wilderness.
	}
}
{
}

\Verse{34}
{
	\hP{בְּמִסְפַּר הַיָּמִים אֲשֶׁר תַּרְתֶּם אֶת הָאָרֶץ אַרְבָּעִים יוֹם יוֹם לַשָּׁנָה יוֹם לַשָּׁנָה תִּשְׂאוּ אֶת עֲוֺנֹתֵיכֶם אַרְבָּעִים שָׁנָה וִידַעְתֶּם אֶת תְּנוּאָתִי׃}
}
{
	\eP{
		For the number of days that you scouted the land, forty
		days, you shall bear your crimes a day for each year,
		forty years, and you shall know my frustration!
	}
}
{
}

\Verse{35}
{
	\hP{אֲנִי יְהֹוָה דִּבַּרְתִּי אִם לֹא זֹאת אֶעֱשֶׂה לְכל הָעֵדָה הָרָעָה הַזֹּאת הַנּוֹעָדִים עָלָי בַּמִּדְבָּר הַזֶּה יִתַּמּוּ וְשָׁם יָמֻתוּ׃}
}
{
	\eP{
		I, YHWH, have spoken: If I shall not do this to all this
		bad congregation who are gathered against me: in this
		wilderness they shall end, and they shall die there!”
	}
}
{
%	\footnote{
%		If I shall not do this. This wording is an oath formulation. It means,
%		“I shall surely do this.”
%	}
}

\Verse{36}
{
	\hP{וְהָאֲנָשִׁים אֲשֶׁר שָׁלַח מֹשֶׁה לָתוּר אֶת הָאָרֶץ וַיָּשֻׁבוּ (וילונו) [וַיַּלִּינוּ] עָלָיו אֶת כּל הָעֵדָה לְהוֹצִיא דִבָּה עַל הָאָרֶץ׃}
}
{
	\eP{
		And the men whom Moses sent to scout the land and came
		back and caused all the congregation to complain against
		him, bringing out a report about the land:
	}
}
{
	Qere Ketiv.
}

\Verse{37}
{
	\hP{וַיָּמֻתוּ הָאֲנָשִׁים מוֹצִאֵי דִבַּת הָאָרֶץ רָעָה בַּמַּגֵּפָה לִפְנֵי יְהֹוָה׃}
}
{
	\eP{
		the men who brought out the bad report of the land died in
		a plague in front of YHWH.
	}
}
{
}

\Verse{38}
{
	\hP{וִיהוֹשֻׁעַ בִּן נוּן וְכָלֵב בֶּן יְפֻנֶּה חָיוּ מִן הָאֲנָשִׁים הָהֵם הַהֹלְכִים לָתוּר אֶת הָאָרֶץ׃}
}
{
	\eP{
		But, out of those men who went to scout the land, Joshua
		son of Nun and Caleb son of Jephunneh lived.  And
		Moses spoke these things to all the children of Israel,
		and the people mourned very much.
	}
}
{
}

\Verse{39}
{
	\hJ{וַיְדַבֵּר מֹשֶׁה אֶת הַדְּבָרִים הָאֵלֶּה אֶל כּל בְּנֵי יִשְׂרָאֵל וַיִּתְאַבְּלוּ הָעָם מְאֹד׃}
}
{
	\eJ{When Moses told these words to all the Israelites, the people mourned greatly.}
}
{
}

\Verse{40}
{
	\hJ{וַיַּשְׁכִּמוּ בַבֹּקֶר וַיַּעֲלוּ אֶל רֹאשׁ הָהָר לֵאמֹר הִנֶּנּוּ וְעָלִינוּ אֶל הַמָּקוֹם אֲשֶׁר אָמַר יְהֹוָה כִּי חָטָאנוּ׃}
}
{
	\eJ{
		And they got up in the morning and went up to the top of
		the mountain, saying, “Here we are, and we’ll go up to the
		place that YHWH said, because we’ve sinned.”
	}
}
{
}

\Verse{41}
{
	\hJ{וַיֹּאמֶר מֹשֶׁה לָמָּה זֶּה אַתֶּם עֹבְרִים אֶת פִּי יְהֹוָה וְהִוא לֹא תִצְלָח׃}
}
{
	\eJ{
		And Moses said, “Why are you violating YHWH’s word? And it
		won’t succeed.
	}
}
{
}

\Verse{42}
{
	\hJ{אַל תַּעֲלוּ כִּי אֵין יְהֹוָה בְּקִרְבְּכֶם וְלֹא תִּנָּגְפוּ לִפְנֵי אֹיְבֵיכֶם׃}
}
{
	\eJ{
		Don’t go up, so you won’t be stricken in front of your
		enemies, because YHWH isn’t among you.
	}
}
{
}

\Verse{43}
{
	\hJ{כִּי הָעֲמָלֵקִי וְהַכְּנַעֲנִי שָׁם לִפְנֵיכֶם וּנְפַלְתֶּם בֶּחָרֶב כִּי עַל כֵּן שַׁבְתֶּם מֵאַחֲרֵי יְהֹוָה וְלֹא יִהְיֶה יְהֹוָה עִמָּכֶם׃}
}
{
	\eJ{
		Because the Amalekite and the Canaanite are there in front
		of you, and you’ll fall by the sword because of the fact
		that you’ve gone back from following YHWH, and YHWH won’t
		be with you.”
	}
}
{
}

\Verse{44}
{
	\hJ{וַיַּעְפִּלוּ לַעֲלוֹת אֶל רֹאשׁ הָהָר וַאֲרוֹן בְּרִית יְהֹוָה וּמֹשֶׁה לֹא מָשׁוּ מִקֶּרֶב הַמַּחֲנֶה׃}
}
{
	\eJ{
		And they acted heedlessly, going up to the top of the
		mountain. And the ark of the covenant of YHWH and Moses
		did not draw away from the camp.
	}
}
{
}

\Verse{45}
{
	\hJ{וַיֵּרֶד הָעֲמָלֵקִי וְהַכְּנַעֲנִי הַיֹּשֵׁב בָּהָר הַהוּא וַיַּכּוּם וַיַּכְּתוּם עַד הַחרְמָה׃}
}
{
	\eJ{
		And the Amalekite and the Canaanite who lived in that
		mountain came down, and they struck them and crushed them
		as far as Hormah.
	}
}
{
}
